\documentclass[12pt,a4paper]{article}

%%%%%%%%%%%%%%%%%%%%%%%%%%%%%%%%%%%%%%%%%%%%%%%%%%
%% DOCUMENT LAYOUT
%%%%%%%%%%%%%%%%%%%%%%%%%%%%%%%%%%%%%%%%%%%%%%%%%%

\usepackage[margin=1.1in]{geometry}
\usepackage{setspace}
\onehalfspacing

\interfootnotelinepenalty=10000


%%%%%%%%%%%%%%%%%%%%%%%%%%%%%%%%%%%%%%%%%%%%%%%%%%
%% FONTS & TYPOGRAPHY
%%%%%%%%%%%%%%%%%%%%%%%%%%%%%%%%%%%%%%%%%%%%%%%%%%

\usepackage[T1]{fontenc}
\usepackage{lmodern}
\usepackage{microtype}


%%%%%%%%%%%%%%%%%%%%%%%%%%%%%%%%%%%%%%%%%%%%%%%%%%
%% MATH & Macros
%%%%%%%%%%%%%%%%%%%%%%%%%%%%%%%%%%%%%%%%%%%%%%%%%%

\usepackage{amsmath,amssymb,amsthm,mathtools}
\usepackage{bbm}
\usepackage{bm}

\DeclareMathOperator{\E}{\mathbb{E}}
\DeclareMathOperator{\Var}{Var}
\DeclareMathOperator{\Cov}{Cov}
\DeclareMathOperator{\Corr}{Corr}
\DeclareMathOperator*{\argmax}{argmax}
\DeclareMathOperator*{\argmin}{argmin}
\DeclareMathOperator{\sign}{sign}
\DeclareMathOperator*{\plim}{plim}
\DeclareMathOperator{\diag}{diag}

\newcommand{\iid}{i.i.d.}
\newcommand{\ind}{\mathbbm{1}}
\newcommand{\indep}{\perp\!\!\!\perp}
\newcommand{\toindistr}{\xrightarrow{d}}

\newcommand{\sumn}{\sum_{i=1}^n}
\newcommand{\wcb}{\overline{\mathrm{Bias}}}
\newcommand{\given}{\mid}

\DeclarePairedDelimiter\abs{\lvert}{\rvert}
\DeclarePairedDelimiter\norm{\lVert}{\rVert}

\newcommand{\defeq}{\vcentcolon=}
\newcommand{\eqdef}{=\vcentcolon}


%%%%%%%%%%%%%%%%%%%%%%%%%%%%%%%%%%%%%%%%%%%%%%%%%%
%% FIGURES & TABLES
%%%%%%%%%%%%%%%%%%%%%%%%%%%%%%%%%%%%%%%%%%%%%%%%%%

\usepackage{graphicx}
\usepackage{tikz}
\usetikzlibrary{decorations.pathreplacing, calc}
\usepackage{booktabs}
\usepackage{subcaption}
\usepackage{siunitx}
\usepackage{float}
\usepackage{placeins}
\usepackage{longtable}

\sisetup{
	group-separator = {,},
	output-decimal-marker = {.},
	group-minimum-digits = 3,
	number-unit-separator = {\,}
}


%%%%%%%%%%%%%%%%%%%%%%%%%%%%%%%%%%%%%%%%%%%%%%%%%%
%% COLORS & LINKS
%%%%%%%%%%%%%%%%%%%%%%%%%%%%%%%%%%%%%%%%%%%%%%%%%%

\usepackage{xcolor}

\usepackage[
colorlinks=true,
linkcolor=blue,
citecolor=blue,
urlcolor=blue
]{hyperref}

\usepackage{cleveref}


%%%%%%%%%%%%%%%%%%%%%%%%%%%%%%%%%%%%%%%%%%%%%%%%%%
%% THEOREMS
%%%%%%%%%%%%%%%%%%%%%%%%%%%%%%%%%%%%%%%%%%%%%%%%%%

\theoremstyle{plain}
\newtheorem{theorem}{Theorem}[section]
\newtheorem{proposition}[theorem]{Proposition}
\newtheorem{lemma}[theorem]{Lemma}
\newtheorem{corollary}[theorem]{Corollary}

\theoremstyle{definition}
\newtheorem{assumption}{Assumption}
\newtheorem{definition}{Definition}

\theoremstyle{remark}
\newtheorem{remark}{Remark}

\crefname{assumption}{assumption}{assumptions}
\Crefname{assumption}{Assumption}{Assumptions}
\Crefname{theorem}{Theorem}{Theorems}

%%%%%%%%%%%%%%%%%%%%%%%%%%%%%%%%%%%%%%%%%%%%%%%%%%
%% NOTES TOGGLE 
%%%%%%%%%%%%%%%%%%%%%%%%%%%%%%%%%%%%%%%%%%%%%%%%%%
\usepackage{todonotes}  % inline colored notes

% --- Global toggle: turn comments on/off ---
\newif\ifcomment
%\commentfalse  % <- switch to \commentfalse to hide comments
\commenttrue

% --- Define comment command with your name ---
\newcommand{\commentT}[1]{%
	\ifcomment
	\todo[inline, color=blue!40]{\textbf{Tobias: }#1}%
	\fi
}
%%%%%%%%%%%%%%%%%%%%%%%%%%%%%%%%%%%%%%%%%%%%%%%%%%
%% USEFUL MACROS
%%%%%%%%%%%%%%%%%%%%%%%%%%%%%%%%%%%%%%%%%%%%%%%%%%




%%%%%%%%%%%%%%%%%%%%%%%%%%%%%%%%%%%%%%%%%%%%%%%%%%
%% BIBLIOGRAPHY (LOAD LAST)
%%%%%%%%%%%%%%%%%%%%%%%%%%%%%%%%%%%%%%%%%%%%%%%%%%

\usepackage[
backend=biber,
style=authoryear,
doi=false,
isbn=false,
url=false
]{biblatex}

\addbibresource{../literature.bib}

%\AtEveryBibitem{\clearfield{note}}
\AtEveryBibitem{%
	\clearfield{note}%
	\clearfield{url}%
	\clearfield{urldate}%
}

\author{Tobias Großbölting}
\title{Meeting Notes --- April 30, 2026}
\begin{document}

\maketitle

\section*{Supervisor Meeting}

\subsection*{General Reception}
Supervisor liked the general idea. Main pushback: ``so what?''\ --- resolved by arguing that the separable class encodes primitive information more tightly than any other class (primitive-consistent, no coupling imported).

\subsection*{Critical Bottleneck: Solvability of the Convex Problem}
Key question: is the modulus/minimax optimization problem still solvable when the constraint set is the separable rectangle, or does the solver rely on the isotropic disk structure?
This is the current bottleneck of the project. If the convex problem remains tractable under the separable class, the main result is essentially complete.

Partial resolution (numerical solvability): for a finite design, the gradient constraint $\nabla f \in \mathcal{C}$ translates to pairwise Lipschitz bounds $f_i - f_j \leq h_{\mathcal{C}}(x_i - x_j)$, which are linear constraints on the function values regardless of $\mathcal{C}$. Only the RHS constants change. The problem has identical structure for both the disk and the rectangle, so numerical solvability is not in question --- this is standard convex optimization.

Open questions:
\begin{itemize}
	\item Can the solution be \emph{characterized analytically} for the rectangle, as AK do for the scalar case (parametric form of least favorable functions, bandwidth parameterization)?
	\item Do AK's regularity conditions (existence/uniqueness of modulus-achieving functions, differentiability of $\omega$) need re-verification for the rectangle?
	\item Does the continuous case (Gaussian white noise, as opposed to finite design) introduce additional difficulties?
	\item Check AK's Supplemental Appendix~E and Kolesár's \texttt{RDHonest} for implicit reliance on scalar/isotropic structure.
\end{itemize}

\subsection*{Wasserstein Distance and Covariate Adjustment}
Supervisor was curious about the 1D Wasserstein distance interpretation in the covariate adjustment problem. He noted that people typically use balancing conditions (balancing of moments).
My counterargument: balancing conditions are imprecise because the worst-case bias function (through the sup) will optimally exploit any slack left by moment-based balancing. The minimax solver already internalizes the Wasserstein distance in the $z$-displacement, just as it internalizes $x$-displacements --- in that sense, it produces an optimal estimator without needing separate balancing conditions.

\subsection*{Action Item: 2D Setup Derivation}
Supervisor wants a clear derivation showing that the 2D setup with treatment rule $T_i = \ind\{X_i \geq c\}$ is equivalent to a univariate RD with covariates. Requirements:
\begin{itemize}
	\item Write in potential outcomes notation.
	\item State a predeterminedness assumption on covariates $Z_i$ and running variable $X_i$.
	\item Show equivalence to univariate RD with covariates explicitly.
\end{itemize}

\subsection*{Action Item: Stefan Wager Poster}
Supervisor mentioned a Stefan Wager poster on optimized RD and covariates. Need to find this poster or contact Stefan Wager directly.

\end{document}
